\documentclass{whutmod}
\usepackage[linesnumbered,ruled,lined]{algorithm2e}
\bibliographystyle{unsrt}
\team{10}
\membera{刘子川}
\joba{编程}
\memberb{程宇}
\jobb{建模}
\memberc{祁成}
\jobc{写作}
\hypersetup{
	colorlinks=true,
	linkcolor=black,citecolor=black
}


\newcommand{\upcite}[1]{\textsuperscript{\cite{#1}}}
%%%%%%%%%%%%%%%%%%%%%%%%%%%%%%%%%题目%%%%%%%%%%%%%%%%%%%%%%%%%%%%%%%%%%%%
\title{基于xxxxxxxx模型}
\tihao{1} 

\begin{document}

	\maketitle
	\thispagestyle{empty}
%%%%%%%%%%%%%%%%%%%%%%%%%%%%%%%%%摘要%%%%%%%%%%%%%%%%%%%%%%%%%%%%%%%%%%%%
	\begin{abstract}
		控制高压油管的压力变化对减小燃油量偏差，提高发动机工作效率具有重要意义。本文建立了基于质量守恒定理的微分方程稳压模型，采用二分法、试探法以及自适应权重的蝙蝠算法对模型进行求解。
		\vspace{6pt}	%空格
	
		针对问题一，建立基于质量守恒定律的燃油流动模型，考察单向阀开启时间对压力稳定性的影响。综合考虑压力与弹性模量、密度之间的关系，提出燃油压力-密度微分方程模型和燃油流动方程。本文采用改进的欧拉方法对燃油压力-密度微分方程求得数值解；利用二分法求解压力分布。综合考虑平均绝对偏差等反映压力稳定程度的统计量，求得直接稳定于100MPa的开启时长为\textbf{0.2955ms} ，在2s、5s内到达并稳定于150MPa时开启时长为\textbf{0.7795ms}、\textbf{0.6734ms}，10s到达并稳定于150MPa的开启时长存在多解。最后对求解结果进行灵敏度分析、误差分析。
		\vspace{6pt}	%空格
	
		针对问题二，建立基于质量守恒定律的泵-管-嘴系统动态稳压模型，将燃油进入和喷出的过程动态化处理。考虑柱塞和针阀升程的动态变动，建立喷油嘴流量方程和质量守恒方程。为提高角速度求解精度，以凸轮转动角度为固定步长，转动时间变动步长，采用试探法粗略搜索与二分法精细搜索的方法求解，求得凸轮最优转动角速度\textbf{0.0283rad/ms（转速270.382转/分钟）}，并得到该角速度下高压油管的密度、压力周期性变化图。对求解结果进行误差分析与灵敏度分析，考察柱塞腔残余容积变动对高压油管压力稳态的影响。
		\vspace{6pt}	%空格
	
		针对问题三，对于增加一个喷油嘴的情况，改变质量守恒方程并沿用问题二的模型调整供、喷油策略，得到最优凸轮转动角速度为\textbf{0.0522rad/ms（498.726转/分钟）}；对于既增加喷油嘴又增加减压阀的情况，建立基于自适应权重的蝙蝠算法的多变量优化模型，以凸轮转动角速度、减压阀开启时长和关闭时长为参数，平均绝对偏差MAD为目标，在泵-管-嘴系统动态稳压模型的基础上进行求解，得到最优参数：\textbf{角速度0.0648 rad/ms（619.109转/分钟）}、减压阀的开启时长\textbf{2.4ms}和减压阀的关闭时长\textbf{97.6ms}。
		\vspace{6pt}	%空格
	
		本文的优点为：1. 采用试探法粗略搜索与二分法精细搜索结合的方法，降低了问题的求解难度。2.以凸轮转动角度为固定步长，对不同角速度按照不同精度的时间步长求解，大大提高了求解的精确度。 3.针对智能算法求解精度方面，采用改进的蝙蝠算法，使速度权重系数自适应调整，兼顾局部搜索与全局搜索能力。
		
		\keywords{
			微分方程\quad
			微分方程\quad	
			微分方程\quad
			微分方程\quad
		}
	\end{abstract}


%%%%%%%%%%%%%%%%%%%%%%%%%%%%%%%%%目录%%%%%%%%%%%%%%%%%%%%%%%%%%%%%%%%%%%%
	\thispagestyle{empty}
	\tableofcontents
	\setcounter{page}{0}                                               
	\newpage	%换页符
	

	
	\section{问题重述}	
		\subsection{问题背景}
	    	新型冠状病毒肺炎（Corona Virus Disease 2019，COVID-19），简称“新冠肺炎”，世界卫生组织命名为“COVID-19”，是指 2019 新型冠状病毒感染导致的肺炎。2020 年 3 月 11 日，世界卫生组织总干事谭德塞宣布，世卫组织认为当前新冠肺炎疫情可被称为全球大流行（pandemic）。目前，COVID-19 疫情仍在世界各地蔓延，已超过 1630 万人感染，65 万余人死亡，给世界各国的经济发展和人民生活带来了极大影响，甚至从一定程度上改变了人类的工作生活方式。
	
		\subsection{问题概述}
		    围绕相关附件和条件要求，定量地研究传染病的传播规律，利用所给（不限于）资料和数据，作出预测并给出控制传染病蔓延的对策建议，具体要求如下：
				 
			
			\textbf{问题一：}
			
			\textbf{问题二：}
			
			\textbf{问题三：}

	
	\section{模型假设}
		\begin{itemize}                                             
		\item [(1)] 
		\item [(2)]
		\item [(3)] 
		\item [(4)] 
		\end{itemize}

		
	\section{符号说明}
		\begin{table}[H]
		\centering
		\setlength{\tabcolsep}{12mm}
		\begin{tabular}{cc}
			\toprule[1.5pt]
			\multicolumn{1}{m{5cm}}{\centering 符号} & \multicolumn{1}{m{5cm}}{\centering 说明} \\
			\midrule[1pt]		
			$P_n$  & 20个站点  \\ 
			$P_n$  & 20个站点  \\ 
		   	$P_n$  & 20个站点  \\ 
			\bottomrule[1.5pt]
		\end{tabular}
		\begin{tablenotes}
		\item 注：表中未说明的符号以首次出现处为准
		\end{tablenotes}
		\end{table}

	\section{问题一模型的建立与求解}
		\subsection{问题描述与分析}
			问题一要求建立至少两个地区的确诊病例数与死亡数的预测模型，并基于模型对这些地区的卫生部门所采取的措施做出评论。收集整理文献可知\upcite{2,3}，灰度预测、时间序列等完全基于数据的预测模型难以更改环境参数，即难以得出卫生部门采取措施对疫情数据的影响。故本节使用可灵活调整参数的动力学模型预测确诊病例数与死亡数的变化趋势。
			
			本文改进了动力学模型中的SEIR模型，基于新冠病毒的特点，添加无症状感染者元素，且使得潜伏者具有传染特性。即总方程组中包含易感者、隔离易感者、潜伏者、隔离潜伏者、感染者以及隔离感染者与无症状感染者等变量。针对方程中的未统计参数，如接触感染率与潜伏者数量等参数，本文使用免疫差分进化算法拟合预测曲线与统计数据以求解该类未知参数。
			
			用改进的动力学模型预测湖北省与美国的确诊病例数与死亡数后，我们通过修改模型参数以探讨当地卫生部门采取的措施对疫情趋势的影响。针对湖北省，我们分别提前和延后$5$天调整模型的隔离比例以分析提前或延后采取严格的隔离措施，对疫情传播所造成的影响。针对美国，我们通过调整日平均接触人数以讨论实施或取消居家令和就地避难令对疫情传播的影响。最后给予部分参数细微震荡，对整体模型进行灵敏度分析。
			
			\begin{figure}[H]
				\centering
				\includegraphics[width=\textwidth]{figures/whut.jpg}
				\caption{问题一思维流程图}\label{lct}
			\end{figure}
			
		\subsection{模型的建立}
			\begin{gather*}
			d(p_i,p_j)=\left | x_i-x_j \right |+\left | y_i-y_j \right | ,
			\end{gather*}
		
		
		\subsection{模型的求解}
		\subsubsection{免疫差分进化算法}
		%\paragraph{初始化}
		本文设计免疫差分进化算法估计微分方程组中的未知参数，定义决策向量为
		\begin{gather}
		X=[x_1,x_2,x_3,x_4],
		\end{gather}
		其中$x_1$、$x_2$、$x_3$和$x_4$分别表示每个患者的日平均接触人数、传染概率、初始潜伏者与患者人数比例以及潜伏者相对于感染者传播能力的比值。将目标函数定义为损失函数如下
		\begin{gather}
		min Loss(X)=\sum_{t=1}^T (|\frac{D_r(t)-D(t)}{D_r(t)}|+|\frac{R_r(t)-R(t)}{R_r(t)}|+|\frac{H_r(t)-H(t)}{H_r(t)}|)
		\end{gather}
		其中$T$表示选取数据的终止节点，即表示选取用于估计参数的数据来自疫情发生的第$1$天到第$T$天。$D_r(t)$、$R_r(t)$与$H_r(t)$分别表示疫情发生后第$t$天的死亡人数、治愈人数和医院患者人数的真实数据；$D(t)$、$R(t)$与$H(t)$分别表示其对应的由SEIR模型。损失函数$Loss$表示预测结果与实际结果间的距离，即$Loss$值越小，预测曲线就与真实曲线越接近。
		\paragraph{种群初始化}
		在解空间中随机产$p$个初始个体$
		X_i(0)=[x_1,x_2,x_3,x_4],(i=1,2,3,\cdots,p).
		$
		其中第$i$个个体的第$j$维取值方式如下
		\begin{gather*}
		x_{i,j}(0)=x_{j,min}+rand(0,1)(x_{j,max}-x_{j,min}),\\i=1,2,3,\cdots,p,j=1,2,3,4
		\end{gather*}
		其中$p$表示种群规模，$x_{j,max}$和$x_{j,min}$分别表示决策变量$X$第$j$维的
		取值范围上界与下界。
		\paragraph{变异}
		在第$g$次迭代中，生成变异个体$H_i(g)$，从种群中随机选取三个个体$X_{p1}(g)$,$X_{p2}(g)$和$X_{p3}(g)$，且$p_1\neq p_2\neq p_3\neq i$，生成的变异向量为
	    \begin{gather}
	    H_i(g)=X_{p1}(g)+F(g)*(X_{p2}(g)-X_{p3}(g))
	    \end{gather}
	    $F(g)\in (0,1)$是每一代中的放缩因子，其服从柯西分部如下
	    \begin{gather*}
	    F(g)=cauchyrnd(uF,0.1)
	    \end{gather*}
	   其中$uF$是$F$的期望值，本文取值为$uF=0.5$。
		\paragraph{交叉}
		对第$g$代种群中第$i$个体进行交叉操作，生成交叉个体$V_i(g)$，具体表达式如下：
		\begin{gather}
		v_{i,j}=\left\{\begin{matrix}h_{i,j}(g),rand(0,1)\leq cr_{i}
		\\ x_{i,j}(g),rand(0,1)>cr_{i}
		\end{matrix}\right.
		\end{gather}
     	其中$cr_{i}\in[0.1,0.6]$是个体$i$的交叉概率，参数$cr_{i}$将进行自适应调整，具体表达式如下：
		\begin{gather}
		cr_{i}=\left\{\begin{matrix}cr_{l}+(cr_{u}-cr_{l})\frac{Loss_{i}-Loss_{min}}{Loss_{max}-Loss_{min}} , Loss_{i}>\overline{Loss}
		\\ cr_{l},Loss_{i}\leqslant  \overline{Loss}
		\end{matrix}\right.
		\end{gather}
	    \paragraph{免疫选择}
	    混合第$g$代的交叉个体$V(g)$与原始个体$X(g)$，得到待选组$\left \{ X '(g+1)\right \}$如下
	    \begin{gather*}
	    X_i '(g+1)=\left\{\begin{matrix}  X_i (g),i\leqslant p
	    \\  V_{i-p} (g),i>p.
	    \end{matrix}\right.
	    \end{gather*}
		个体 $X_a '(g+1)$和$X_b '(g+1)$的亲和度$S_{a,b}$可表示为
		\begin{gather}
		S_{a,b}=\sqrt{\sum _{i=1}^4( \frac{x_{i,a}-x_{i,b}}{x_{i,max}-x_{i,min}})^2}
		\end{gather}
		$S_{a,b}$为$X_a '(g+1)$和$X_b '(g+1)$的归一化距离，表示个体$X_a '(g+1)$和$X_b '(g+1)$的相似性。定义个体$X_i '(g+1)$的抗体浓度为$C_{i}$,即
		\begin{gather}
		C_{i}=\frac{1}{2p}\sum _{j=1}^{2p} N_{i,j},\\
		N_{i,j}=\left\{\begin{matrix}1,S_{i,j}\geqslant \mu 
		\\ 0,S_{i,j}< \mu 
		\end{matrix}\right.
		\end{gather}
		$\mu(\mu\in[0,1])$为相似度阈值，即当个体$i$和$j$的亲和度$S_{i,j}\geqslant \mu$时认为个体$i$和$j$为相似个体。$C_{i}$即为$\left \{ X '(g+1)\right \}$中$X_i '(g+1)$的相似个体所占比例，$C_{i}$越大即表示$X_i '(g+1)$所在区域的个体密度越大。我们优先将损失函数$Loss$值最优的前$\sigma$个解放入下一代个体$\left \{ X(g+1)\right \}$中以防止最优解丢失。再计算剩余个体的复合适应度函数，即个体$i$的复合适应度函数可表示为
	    \begin{gather}
	  min F(X_i '(g+1))=\frac{Loss(X_i '(g+1))-Loss_{min}}{Loss_{max}-loss_{min}}+C_{i}
		\end{gather}
		即选取复合适应度函数$F$较优的剩余$p-\sigma$个个体放入下一代个体$\left \{ X(g+1)\right \}$中。重复迭代上述算法$G$次后终止算法并输出最优参数集$X_{best}$。
		
		
		
%				\begin{algorithm}[H]
%			 	\caption{Procedure of Apriori}  
%			 	\LinesNumbered  
%			 	\setstretch{0.9}   %设置表的行间距
%			 	\KwIn{item data base: $D$\newline
%			 		minimum Support threshold: $Sup_{min}$\newline
%			 		minimum Confidence threshold: $Conf_{min}$
%			 	}
%			 	\KwOut{frequent item sets $F$}  
%			 	\textbf{Initialize} \newline
%			 	iteration $t\leftarrow 1$ \newline
%			 	The candidate FIS:$C_{t}=\varnothing$ \newline
%			 	The length of FIS:$length=1$ \newline
%			 	\For{i=1 to sizeof(D)}
%			 	{$I_{i}$=D(i)\newline
%			 		n=sizeof($I_{i}$)\newline
%			 		\For{j=1 to n}{
%			 			\If{$I_{i}(j)\notin C_{t}$ }
%			 			{$C_{t}=C_{t}\cup I_{i}(j) $}
%			 		}
%			 	}
%			 	$F_{t}=\left \{ f|f\in C_{t},Sup(f)>Sup_{min}\right \}$\newline
%			 	\While{$F\neq \varnothing$}
%			 	{ t=t+1\newline 
%			 		length=length+1\newline	
%			 		$C_{t}\leftarrow $ all candidate of FIS in $F_{t-1}$\newline
%			 		$F_{t}=\left \{ f|f\in C_{t},(Sup(f)>Sup_{min})\bigcap (Comf(f)>Conf_{min}) \right\}$\newline
%			 	}	
%			 	\Return{$F_{t-1}$} 
%			\end{algorithm} 
		
        \subsection{实验结果及分析}
  
	\section{问题二模型的建立与求解}
		\subsection{问题描述与分析}
			问题二要求

    		其思维流程图如图~\ref{lssssct}~所示：

			\begin{figure}[H]
				\centering
				\includegraphics[width=\textwidth]{figures/whut.jpg}
				\caption{问题二思维流程图}\label{lssssct}
			\end{figure}

		\subsection{模型的建立}
		
		\subsection{模型的求解}
	
		
             
        \subsection{实验结果及分析}
        
			结果如下表\ref{zhuanssssasgzai}所示：
			\begin{table}[H]
			\setstretch{1.4}  %设置表的行间距
			\centering		
			\caption{xxxxxxxxxxxxxxxxxxxxx}\label{biao1}
			\begin{tabular}{cc}
			\toprule[2pt]
				\multicolumn{1}{m{5cm}}{\centering xxxxxxx}
				& \multicolumn{1}{m{5cm}}{\centering xxxxxxx}
				\\
				\midrule[1pt]
				xxxxxxx &   909.80\\ 
				xxxxxxx & 	852.60\\ 
			\bottomrule[2pt]	
			\end{tabular}
			\end{table}
  
  			由表\ref{biao1}可知

			其各个小车的运输细节图下图所示：
			\begin{figure}[H]
				\centering
				\subfigure{\includegraphics[height=8cm,width=7.5cm]{figures/whut.jpg}}
				\subfigure{\includegraphics[height=8cm,width=7.5cm]{figures/whut.jpg}}
			\end{figure}	
			\begin{figure}[H]	
				\centering
				\subfigure{\includegraphics[height=8cm,width=7.5cm]{figures/whut.jpg}}
				\subfigure{\includegraphics[height=8cm,width=7.5cm]{figures/whut.jpg}}
				\caption{xxxxxxxxxxxxxxxxxxxxxxxxx}
				\label{fisg}
			\end{figure}

    \section{问题三模型的建立与求解}
  		\subsection{结果分析}
  
  	\section{灵敏度分析}
 
  	\section{模型的评价}
		\subsection{模型的优点}
			\begin{itemize}                                             
			\item [(1)]
			\item [(2)] 	
			\end{itemize}
		\subsection{模型的缺点}

  		\subsection{模型改进}

  
  
 
	\newpage	%换页符
	%%参考文献
	%\begin{thebibliography}{9}%宽度9
	% \setlength{\itemsep}{-2mm}
	\nocite{*}		%排版未引用的参考文献
	\begin{thebibliography}{9}%宽度9
		\bibitem{1}张斯嘉, 郭建胜, 钟夫, 等. 基于蝙蝠算法的多目标战备物资调运决策优化[J]. 火力与指挥控制, 2016, 41(1): 58-61.
	
	\end{thebibliography}

	\newpage
	%附录
	\appendix %%附录
	\section{数据可视化的实现}
		\subsection*{第一问画图--python源代码}
			\begin{lstlisting}[language=python]
			
			\end{lstlisting}
			
		\subsection*{第二问画图--python源代码}
			\lstinputlisting[language={python},numbers=left,numberstyle=\tiny,
			rulesepcolor=\color{red!20!green!20!blue!20},  
			keywordstyle=\color{blue!70!black},  
			commentstyle=\color{blue!90!},  
			basicstyle=\ttfamily] {./code/demo.py}

\end{document}